\documentclass[12pt]{article}

\usepackage[a4paper,margin=1in]{geometry}
\usepackage[T1]{fontenc}
\usepackage[utf8]{inputenc}
\usepackage{lmodern}
\usepackage{amsmath,amssymb}
\usepackage{booktabs}
\usepackage{longtable}
\usepackage{array}
\usepackage{tabularx}
\usepackage{enumitem}
\usepackage{hyperref}
\usepackage{setspace}
\usepackage{ragged2e}

\setstretch{1.08}
\setlength{\parskip}{0.5em}
\setlength{\parindent}{0pt}

\hypersetup{
	colorlinks=true,
	linkcolor=black,
	citecolor=black,
	urlcolor=blue,
	pdftitle={Collected Propositions and Theorems for the General Theory of Cognitive Structuring},
	pdfauthor={Kostiantyn Osmolovskyi}
}

\setlist[itemize]{topsep=3pt,itemsep=2pt,parsep=0pt}
\setlist[enumerate]{topsep=3pt,itemsep=2pt,parsep=0pt}


\title{Collected Propositions and Theorems \\
	\small for the General Theory of Cognitive Structuring}

\author{
	Kostiantyn Osmolovskyi\\
	\small Independent Researcher, Ukraine\\
	\small ORCID: 0009-0006-3144-7237\\
	\small \texttt{constantinosmol@gmail.com}
}

\date{}

% ---------------------------------------------------------------------------
% Part of a series: General Theory of Cognitive Structuring (GTCS)
% Autor: Kostiantyn Osmolovskyi (ORCID: 0009-0006-3144-7237)
% License: Creative Commons Attribution 4.0 International (CC BY 4.0)
% Repository: https://github.com/k-osmolovskyi/k-osmolovskyi.github.io
% DOI: https://doi.org/10.5281/zenodo.19705678
% ---------------------------------------------------------------------------

\begin{document}
	\maketitle

\begin{center}
	\small
	Verification Report No. 2026-2\\
	Version 1.0
\end{center}

\vspace{0.5em}	
	
	\section{Collected geometric propositions and theorems}
	
	\subsection*{G1.01. Coherence as distance to the stability region}
	\textbf{Type:} Definition-formal proposition \\
	\textbf{Statement:} Let $X$ be the configuration space of processing states and let
	$U \subseteq X$ denote the region of structurally stable configurations determined by the
	architecture. Then coherence is defined as the geometric distance between the current
	configuration $x_t$ and the stability region:
	\[
	C(x_t) = \mathrm{dist}(x_t,U).
	\]
	In later invariant-grounded form, this may be written as
	\[
	C_I(x_t) = \mathrm{dist}(x_t,U_X(I_t)).
	\]
	\textbf{Dependencies:} configuration space; current configuration; stability region; distance structure. \\
	\textbf{Assumptions:} existence of an admissible metric or distance relation on the configuration space. \\
	\textbf{Main source:} \emph{Coherence Evaluation in Cognitive Systems}. \\
	\textbf{Proof status:} P. \\
	\textbf{Note:} This result is foundational for the geometric branch. The later notation
	$U_X(I_t)$ strengthens the architectural grounding of the earlier region $U$ without changing
	the role of coherence as a geometric quantity.
	
	\subsection*{G1.02. Coherence is distinct from structural tension}
	\textbf{Type:} Proposition \\
	\textbf{Statement:} Coherence and structural tension are not identical quantities. Coherence
	measures the geometric deviation of the current configuration from the stability region,
	whereas structural tension measures the instantaneous regulatory burden required to sustain
	processing under interacting constraints. Accordingly, structural tension may function as a
	local correlate of deviation from stability without replacing the geometric notion of coherence. \\
	\textbf{Dependencies:} coherence as distance to the stability region; tension decomposition
	$T_t = K_t + C_t$. \\
	\textbf{Assumptions:} none beyond the architectural distinction between global geometric relation
	and local regulatory burden. \\
	\textbf{Main source:} \emph{Coherence Evaluation in Cognitive Systems}. \\
	\textbf{Proof status:} C. \\
	\textbf{Note:} This is a normalized proposition-level formulation of a distinction argued
	explicitly in the coherence paper. It should not be read as claiming that no local relation
	between coherence and tension can be studied, only that they are not the same formal object.
	
	\section{Collected regulatory and overload propositions and theorems}
	
	\subsection*{R2.01. Structural admissibility depends on the joint regulatory state}
	\textbf{Type:} Proposition \\
	\textbf{Statement:} Structural admissibility of updating is not determined by instantaneous
	tension alone. It is determined by the joint regulatory state including at least current
	tension and accumulated overload memory. Equivalently, admissibility is defined over a
	two-dimensional regulatory phase space rather than over present tension in isolation. \\
	\textbf{Dependencies:} structural tension; overload memory; admissibility operator
	$A(T,L,R,M,\bar h)$. \\
	\textbf{Assumptions:} limited resources; overload memory; admissibility operator defined over
	joint regulatory variables. \\
	\textbf{Main source:} \emph{Structural Admissibility in Cognitive Systems}. \\
	\textbf{Proof status:} P. \\
	\textbf{Note:} This is one of the central results of the regulatory branch and grounds the
	phase-space treatment used in later papers.
	
	\subsection*{R2.02. Overload memory is not reducible to instantaneous tension}
	\textbf{Type:} Proposition \\
	\textbf{Statement:} Overload memory is a trajectory-dependent state variable and cannot, in
	general, be reduced to a function of instantaneous tension alone. That is, there is no
	function $f$ such that
	\[
	L_t = f(T_t)
	\]
	for all admissible trajectories of the architecture. \\
	\textbf{Dependencies:} overload excess $u_t$; overload update rule; retained state variable
	$L_t$. \\
	\textbf{Assumptions:} nontrivial temporal accumulation of excess load; admissible trajectories
	with distinct overload histories. \\
	\textbf{Main source:} \emph{Overload Formation in Cognitive Processing}. \\
	\textbf{Proof status:} P. \\
	\textbf{Note:} This result establishes overload memory as a genuine additional state dimension
	of regulation.
	
	\subsection*{R2.03. Regulatory state is trajectory-dependent}
	\textbf{Type:} Proposition \\
	\textbf{Statement:} The regulatory state of a cognitive system cannot be determined from
	instantaneous tension alone. If two trajectories satisfy
	\[
	T_t^{(1)} = T_t^{(2)}
	\]
	at some time $t$ but differ in overload memory,
	\[
	L_t^{(1)} \neq L_t^{(2)},
	\]
	then their admissibility conditions may differ:
	\[
	A\!\left(T_t^{(1)},L_t^{(1)},\dots\right)
	\neq
	A\!\left(T_t^{(2)},L_t^{(2)},\dots\right).
	\]
	\textbf{Dependencies:} overload memory; admissibility operator; joint-state admissibility. \\
	\textbf{Assumptions:} admissibility depends jointly on instantaneous tension and overload memory. \\
	\textbf{Main source:} \emph{Trajectory-Dependent Regulation in Cognitive Systems}. \\
	\textbf{Proof status:} P. \\
	\textbf{Note:} This proposition makes explicit the path dependence implicit in overload-based
	regulation.
	
	\subsection*{R2.04. Overload deforms the admissibility boundary and introduces hysteresis}
	\textbf{Type:} Proposition \\
	\textbf{Statement:} If admissibility in regulatory phase space is described by a critical
	boundary
	\[
	T_{\mathrm{crit}}(L),
	\]
	and if this boundary decreases monotonically with overload memory, then increasing overload
	reduces the maximal instantaneous tension that can be tolerated without structural updating.
	As a consequence, identical instantaneous tensions may lie inside or outside the stability
	region depending on overload history, yielding regulatory hysteresis. \\
	\textbf{Dependencies:} stability region in $(T,L)$-space; critical tension function
	$T_{\mathrm{crit}}(L)$; overload memory dynamics. \\
	\textbf{Assumptions:} monotonic boundary deformation with respect to overload memory. \\
	\textbf{Main source:} \emph{Trajectory-Dependent Regulation in Cognitive Systems}. \\
	\textbf{Proof status:} P/S. \\
	\textbf{Note:} The geometric claim is explicit and central; the exact strength of the result
	depends on the boundary-deformation assumption adopted in the trajectory paper.
	
	\subsection*{R2.05. The regulatory phase space is minimally two-dimensional}
	\textbf{Type:} Proposition \\
	\textbf{Statement:} The minimal regulatory state sufficient to characterize admissibility
	dynamics includes at least two variables: instantaneous tension $T$ and overload memory $L$.
	Accordingly, the minimal regulatory phase space is
	\[
	\Omega = \{(T,L)\mid T \ge 0,\; L \ge 0\}.
	\]
	\textbf{Dependencies:} structural tension; overload memory; joint-state admissibility. \\
	\textbf{Assumptions:} overload is retained as a genuine state component rather than collapsed
	into present pressure. \\
	\textbf{Main source:} \emph{Trajectory-Dependent Regulation in Cognitive Systems}, grounded in
	\emph{Structural Admissibility in Cognitive Systems} and \emph{Overload Formation in Cognitive Processing}. \\
	\textbf{Proof status:} C. \\
	\textbf{Note:} This is a canonical collected proposition: it is strongly supported by the
	regulatory branch even though the two-dimensionality is distributed across more than one paper.
	
	\section{Collected invariant and compression propositions and theorems}
	
	\subsection*{R3.01. Invariant-induced stability geometry}
	\textbf{Type:} Proposition/Theorem \\
	\textbf{Statement:} Let $I_t=\{I_1,\dots,I_n\}$ be the current invariant structure of the
	architecture, and let each invariant $I_i$ admit a local admissibility profile
	\[
	a_i(x)\in[0,1]
	\]
	over the configuration space $X$. Then the invariant structure induces a region of admissible
	or stable configurations, written in later grounded form as $U_X(I_t)$, relative to which
	geometric coherence may be defined. In this sense, the stability geometry of the system is not
	primitive but structurally induced by the current invariant organization. \\
	\textbf{Dependencies:} invariant; invariant structure; local admissibility profiles; thresholded
	or admissible configuration conditions. \\
	\textbf{Assumptions:} sufficient regularity of admissibility profiles and nonempty admissible
	configuration set. \\
	\textbf{Main source:} \emph{Invariants in Cognitive Architectures}. \\
	\textbf{Proof status:} P. \\
	\textbf{Note:} This result provides the architectural grounding of the earlier stability-region
	notation used in the coherence paper.
	
	\subsection*{R3.02. Global conflict is structured by invariant interaction architecture}
	\textbf{Type:} Proposition \\
	\textbf{Statement:} Global structural conflict is determined not by isolated local deviations
	alone but by the interaction architecture of invariants. More precisely, if pairwise
	incompatibilities are defined over the interaction graph
	\[
	G_t=(V_t,E_t), \qquad V_t=I_t,
	\]
	and aggregated by a monotone conflict functional, then the resulting incompatibility load
	$K_t$ reflects the structural organization of constraint interactions rather than merely the
	presence of independent local tensions. \\
	\textbf{Dependencies:} invariant structure; interaction graph; pairwise incompatibility
	functions; conflict aggregation. \\
	\textbf{Assumptions:} edge locality, monotone aggregation, and admissible conflict
	construction. \\
	\textbf{Main source:} \emph{Invariants in Cognitive Architectures}. \\
	\textbf{Proof status:} P. \\
	\textbf{Note:} This proposition is central for understanding why tension and overload are
	architectural rather than merely pointwise quantities.
	
	\subsection*{R3.03. Compression is not equivalent to simplification}
	\textbf{Type:} Proposition \\
	\textbf{Statement:} Structural compression and structural simplification are not equivalent.
	An admissible update may reduce the cardinality or apparent complexity of invariant structure
	without reducing expected overload. Therefore, reduction in structural size is not sufficient
	to characterize stabilizing architectural change. \\
	\textbf{Dependencies:} admissible structural updating; invariant structure; expected overload
	criterion. \\
	\textbf{Assumptions:} compression evaluated with respect to long-run regulatory effect under
	processing dynamics. \\
	\textbf{Main source:} \emph{Structural Compression in Cognitive Architectures}. \\
	\textbf{Proof status:} P. \\
	\textbf{Note:} This result blocks the collapse of all admissible structural reduction into one
	class of beneficial transformation.
	
	\subsection*{R3.04. Compression generates a new invariant}
	\textbf{Type:} Proposition \\
	\textbf{Statement:} Let
	\[
	I_{t+1}=(I_t\setminus S)\cup\{F(S)\},
	\]
	where $S\subset I_t$ and the update is a structural compression. Then the compressed
	representation $F(S)$ enters the updated architecture as a new invariant. In this sense,
	compression is generative of invariant structure and is not merely eliminative. \\
	\textbf{Dependencies:} compression operation; invariant definition; admissible structural update. \\
	\textbf{Assumptions:} compressed representation participates in later processing and is
	preserved under ordinary processing at the updated level. \\
	\textbf{Main source:} \emph{Structural Compression in Cognitive Architectures}. \\
	\textbf{Proof status:} P. \\
	\textbf{Note:} This proposition is one of the main bridges between the invariant branch and the
	architectural evolution branch.
	
	\subsection*{R3.05. Compression can preserve or transform admissible geometry}
	\textbf{Type:} Proposition \\
	\textbf{Statement:} Structural compression may either preserve the equivalence class of the
	admissible stability geometry or alter it. Accordingly, not every compression yields a level
	transition. Some compressions are level-preserving, while others are level-forming because they
	change the geometry of admissibility in a non-equivalent way. \\
	\textbf{Dependencies:} compression operation; admissible geometry; architectural level. \\
	\textbf{Assumptions:} geometric comparison of pre-compression and post-compression stability
	regions under admissible reparameterization. \\
	\textbf{Main source:} \emph{Structural Compression in Cognitive Architectures}. \\
	\textbf{Proof status:} P/S. \\
	\textbf{Note:} The distinction is clearly established in the compression paper; exact theorem
	strength may later benefit from more explicit normalized extraction.
	
	\subsection*{R3.06. Level transition is a non-equivalent change in admissible geometry}
	\textbf{Type:} Proposition \\
	\textbf{Statement:} A transition between architectural levels occurs if and only if the
	stability geometry induced by the updated architecture is not equivalent to the prior one under
	the admissible equivalence relation used for level definition. Thus, level transition is not a
	matter of more information, more intensity, or greater complexity alone, but of non-equivalent
	geometric reorganization. \\
	\textbf{Dependencies:} architectural level; stability geometry; geometric equivalence relation. \\
	\textbf{Assumptions:} level defined by equivalence class of admissible geometry. \\
	\textbf{Main source:} \emph{Structural Admissibility in Cognitive Systems}, strengthened in
	\emph{Structural Compression in Cognitive Architectures}. \\
	\textbf{Proof status:} P/C. \\
	\textbf{Note:} This result should be read as a collected canonical proposition whose grounding
	begins in the structural admissibility paper and becomes operationally sharper in the
	compression paper.
	
	\subsection*{R3.07. Compression improves effective long-run regulatory stability}
	\textbf{Type:} Proposition \\
	\textbf{Statement:} If an admissible structural update qualifies as a compression in the strict
	sense, then it reduces expected overload and thereby improves the effective long-run regulatory
	stability of the architecture. This improvement is long-run and regulatory rather than merely
	local or instantaneous. \\
	\textbf{Dependencies:} compression criterion; expected overload; overload-sensitive regulatory
	evaluation. \\
	\textbf{Assumptions:} compression preserves admissible processing and reduces expected overload
	under the induced processing distribution. \\
	\textbf{Main source:} \emph{Structural Compression in Cognitive Architectures}. \\
	\textbf{Proof status:} P/S. \\
	\textbf{Note:} The paper states this result in effective long-run form rather than as a
	pointwise theorem about all admissibility boundaries; that narrower scope should be preserved in
	canonical normalization.
	
	\section{Collected representation and accessibility propositions and theorems}
	
	\subsection*{R4.01. Coherence representation is an internal regulatory variable rather than a reconstruction of geometric coherence}
	\textbf{Type:} Proposition \\
	\textbf{Statement:} Coherence representation is an internal regulatory variable constructed
	from structurally relevant observable signals and used for regulation under bounded access.
	It does not reconstruct the full geometric coherence relation
	\[
	C_I(x_t)=\mathrm{dist}(x_t,U_X(I_t)),
	\]
	and should not be identified with geometric coherence itself. \\
	\textbf{Dependencies:} geometric coherence; observable structural signals; internal
	representation map. \\
	\textbf{Assumptions:} regulation proceeds under partial observability of global architecture. \\
	\textbf{Main source:} \emph{Emergence of Coherence Representation}. \\
	\textbf{Proof status:} P. \\
	\textbf{Note:} This is the central entry proposition of the representational branch and fixes
	the non-reconstructive status of coherence representation.
	
	\subsection*{R4.02. Coherence representation preserves regulatory order without requiring metric preservation}
	\textbf{Type:} Proposition \\
	\textbf{Statement:} For regulatory purposes, coherence representation need only preserve the
	ordering of states with respect to regulatory significance. It need not preserve the full metric
	structure of geometric coherence. Accordingly, two states may be correctly ordered for
	regulation even if their exact geometric distances are not internally reconstructed. \\
	\textbf{Dependencies:} coherence representation; regulatory significance ordering; geometric
	coherence. \\
	\textbf{Assumptions:} bounded observability and order-sensitive regulation. \\
	\textbf{Main source:} \emph{Emergence of Coherence Representation}, strengthened in
	\emph{Coherence Representation in Multi-System Regulation}. \\
	\textbf{Proof status:} P/C. \\
	\textbf{Note:} This is a canonical collected proposition spanning the single-system and
	multi-system representation papers.
	
	\subsection*{R4.03. Pre-symbolic admissibility is a distinct prior filtering layer}
	\textbf{Type:} Proposition-like definitional result \\
	\textbf{Statement:} Pre-symbolic admissibility is a distinct regulatory layer that operates on
	the pre-representational discrepancy or signal domain prior to stable representation and prior
	to structural updating. It must therefore be kept distinct from both coherence representation
	and the admissibility operator for structural updating. \\
	\textbf{Dependencies:} pre-symbolic signal domain; regulatory layering; distinction between
	representation and updating. \\
	\textbf{Assumptions:} none beyond the layered architecture of access and regulation. \\
	\textbf{Main source:} \emph{Pre-Symbolic Admissibility in Cognitive Systems}. \\
	\textbf{Proof status:} D/C. \\
	\textbf{Note:} Although this is fundamentally a definitional-layer result, it is included here
	because it functions as a canonical proposition-level separation of operators:
	\[
	\Pi_t^{ps} \neq \widehat{C}_t \neq A(\cdot).
	\]
	
	\subsection*{R4.04. Pre-symbolic admissibility determines the admissible discrepancy domain}
	\textbf{Type:} Proposition \\
	\textbf{Statement:} Let $Z_t$ denote the pre-symbolic discrepancy domain and let
	\[
	\Pi_t^{ps}: Z_t \to \{0,1\}^m
	\]
	denote the pre-symbolic admissibility operator. Then the admissible discrepancy domain
	\[
	\mathcal{D}_t := Z_t^{adm}
	\]
	is determined by the discrepancies admitted by $\Pi_t^{ps}$. Only those discrepancies that
	enter $\mathcal{D}_t$ may contribute to later observable signal formation and subsequent regulatory
	representation. \\
	\textbf{Dependencies:} pre-symbolic signal domain; pre-symbolic admissibility operator. \\
	\textbf{Assumptions:} admitted discrepancies are the only discrepancies propagated beyond the
	pre-symbolic stage. \\
	\textbf{Main source:} \emph{Pre-Symbolic Admissibility in Cognitive Systems}. \\
	\textbf{Proof status:} P/S. \\
	\textbf{Note:} This is the main domain-forming proposition of the access branch.
	
	\subsection*{R4.05. Restricted accessibility of coherence}
	\textbf{Type:} Proposition \\
	\textbf{Statement:} Geometric coherence may exist without becoming regulatorily accessible.
	More precisely, if the discrepancies corresponding to a geometric deviation from the stability
	region are not admitted into the admissible discrepancy domain $\mathcal{D}_t$, then no corresponding
	regulatory signal is generated. Thus coherence, as a geometric property of system state, and
	coherence, as a regulatorily accessible quantity, must be distinguished. \\
	\textbf{Dependencies:} geometric coherence; admissible discrepancy domain; regulatory
	accessibility; coherence representation. \\
	\textbf{Assumptions:} pre-symbolic admissibility acts prior to signal-set construction and
	representation formation. \\
	\textbf{Main source:} \emph{Restricted Accessibility of Coherence in Cognitive Systems}. \\
	\textbf{Proof status:} P. \\
	\textbf{Note:} This is the principal bridge proposition between the geometric branch and the
	access branch.
	
	\subsection*{R4.06. Non-injectivity of regulatory accessibility}
	\textbf{Type:} Proposition \\
	\textbf{Statement:} Distinct configurations may become indistinguishable at the level of
	regulatorily accessible signal sets. In particular, there may exist
	\[
	x_t \neq x_t'
	\]
	such that the configurations differ geometrically while yielding identical admitted observable
	signal sets and hence identical or indistinguishable accessible regulatory representations. \\
	\textbf{Dependencies:} restricted accessibility of coherence; admissible discrepancy domain;
	observable signal-set construction. \\
	\textbf{Assumptions:} differences between configurations may fail to enter the admissible
	discrepancy domain. \\
	\textbf{Main source:} \emph{Restricted Accessibility of Coherence in Cognitive Systems}. \\
	\textbf{Proof status:} P/S. \\
	\textbf{Note:} This proposition makes explicit that regulatorily accessible structure is not an
	injective image of full configuration structure.
	
	\subsection*{R4.07. Restricted accessibility is structural, not merely representationally approximate}
	\textbf{Type:} Proposition \\
	\textbf{Statement:} The indistinguishability induced by restricted accessibility is not, in
	general, a consequence of noise, approximation error, or limited representational resolution.
	It is a structural consequence of admissibility constraints operating prior to stable
	representation. \\
	\textbf{Dependencies:} restricted accessibility of coherence; non-injectivity of regulatory
	accessibility; pre-symbolic admissibility. \\
	\textbf{Assumptions:} pre-symbolic admissibility determines signal availability rather than
	merely filtering already formed stable representations. \\
	\textbf{Main source:} \emph{Restricted Accessibility of Coherence in Cognitive Systems}. \\
	\textbf{Proof status:} C. \\
	\textbf{Note:} This is a canonical normalized proposition drawn from the interpretive and
	architectural consequences sections of the restricted accessibility paper.
	
	\subsection*{R4.08. Repeated admissible propagation induces propagation bias}
	\textbf{Type:} Proposition \\
	\textbf{Statement:} If certain classes of discrepancies are repeatedly admitted at the
	pre-symbolic stage while others are repeatedly excluded, the resulting propagation profile is
	not neutral. Over time, this selective propagation induces stable biases in the discrepancy
	classes that are available to downstream regulation and representation. \\
	\textbf{Dependencies:} pre-symbolic admissibility; admissible discrepancy domain; repeated
	propagation under ongoing processing. \\
	\textbf{Assumptions:} temporally repeated admissibility filtering with nonuniform admission
	patterns. \\
	\textbf{Main source:} \emph{Pre-Symbolic Admissibility in Cognitive Systems}. \\
	\textbf{Proof status:} S. \\
	\textbf{Note:} This proposition is central for the dynamic significance of the pre-symbolic
	branch, but its canonical extraction may later benefit from fuller proof normalization.
	
	\subsection*{R4.09. Repeated admissible propagation may stabilize proto-structural channels}
	\textbf{Type:} Proposition \\
	\textbf{Statement:} Under repeated admissible propagation, certain discrepancy-transmission
	patterns may become sufficiently stable to function as proto-structural channels for later
	regulation and, potentially, for later representational consolidation. \\
	\textbf{Dependencies:} pre-symbolic admissibility; propagation bias; repeated admissible signal
	selection. \\
	\textbf{Assumptions:} persistence of selective propagation dynamics under bounded perturbation. \\
	\textbf{Main source:} \emph{Pre-Symbolic Admissibility in Cognitive Systems}. \\
	\textbf{Proof status:} S/O. \\
	\textbf{Note:} This result is important for the developmental meaning of the access layer, but
	should currently be treated as a sketch-level or pending-consolidation proposition until a more
	fully normalized proof layer is assembled.
	
	\section{Collected identity and multi-system propositions and theorems}
	
	\subsection*{R5.01. Identity is derived as a regulatory attractor rather than introduced as a primitive self-representational object}
	\textbf{Type:} Proposition \\
	\textbf{Statement:} Within the present framework, identity is not posited as a primitive
	self-representational substance or as a basic symbolic object. Rather, identity is derived as a
	regulatory consequence of bounded dynamics in which trajectories of the system tend to remain
	within, return to, or concentrate around regions of comparatively low overload and sustained
	structural admissibility. \\
	\textbf{Dependencies:} overload memory; admissibility boundary; trajectory-dependent regulation;
	bounded regulatory dynamics. \\
	\textbf{Assumptions:} existence of a long-run regulatory regime in which recurrent low-overload
	regions are dynamically meaningful. \\
	\textbf{Main source:} \emph{Identity as a Regulatory Attractor}. \\
	\textbf{Proof status:} S/C. \\
	\textbf{Note:} This proposition should be read as the canonical entry claim of the identity
	branch. Its role is to block the reduction of identity to narrative or symbolic self-description.
	
	\subsection*{R5.02. Low-overload concentration regions may function as identity-cores}
	\textbf{Type:} Proposition/Theorem \\
	\textbf{Statement:} Under suitable long-run regulatory conditions, trajectories of the system
	may concentrate near regions of comparatively low overload and persistent admissibility. Such
	regions may function as identity-cores in the sense that they support recurrent organization
	and relative structural persistence across time. \\
	\textbf{Dependencies:} overload memory dynamics; admissibility structure; trajectory dependence;
	regulatory phase space. \\
	\textbf{Assumptions:} bounded regulatory dynamics; long-run concentration or stationarity-type
	conditions local to the identity branch. \\
	\textbf{Main source:} \emph{Identity as a Regulatory Attractor}. \\
	\textbf{Proof status:} S/P. \\
	\textbf{Note:} The exact theorem strength depends on the long-run assumptions adopted in the
	identity paper. In canonical use, the result should remain tied to those local assumptions.
	
	\subsection*{R5.03. Identity persistence is regulatory rather than merely representational}
	\textbf{Type:} Proposition \\
	\textbf{Statement:} Persistence of identity is not explained solely by persistence of
	self-representation or symbolic self-description. In the present framework, identity persistence
	is grounded in recurrent regulatory organization, especially in the tendency of the system to
	remain near or return toward regions of structurally sustainable low-overload organization. \\
	\textbf{Dependencies:} identity as regulatory attractor; low-overload concentration regions;
	bounded regulation. \\
	\textbf{Assumptions:} identity interpreted architecturally rather than semantically. \\
	\textbf{Main source:} \emph{Identity as a Regulatory Attractor}. \\
	\textbf{Proof status:} C. \\
	\textbf{Note:} This is a canonical normalized proposition drawn from the architectural reading
	of the identity branch rather than from a single theorem statement.
	
	\subsection*{R5.04. Inter-system comparability does not require shared metric coherence or shared configuration space}
	\textbf{Type:} Proposition \\
	\textbf{Statement:} Two systems may become comparable for regulatory purposes without sharing a
	common metric geometry of coherence and without sharing a common configuration space. What is
	required is not metric identity of state space, but a transferable regulatory relation between
	their internally available coherence-related variables. \\
	\textbf{Dependencies:} coherence representation; non-metric regulatory ordering; multi-system
	setting. \\
	\textbf{Assumptions:} systems possess coherence-related internal regulatory variables adequate
	for ordered comparison. \\
	\textbf{Main source:} \emph{Coherence Representation in Multi-System Regulation}. \\
	\textbf{Proof status:} P/C. \\
	\textbf{Note:} This proposition is the entry point of the multi-system branch and prevents false
	assumptions of shared geometry.
	
	\subsection*{R5.05. Minimal transferable structure between systems is ordinal rather than metric}
	\textbf{Type:} Proposition \\
	\textbf{Statement:} The minimal structure that must be preserved for inter-system regulatory
	transfer is not full metric equality of coherence representation, but preservation of
	regulatorily relevant order. Accordingly, if there exists an order-preserving map
	\[
	\phi_{AB}
	\]
	between the images of coherence representations in systems $A$ and $B$, then cross-system
	regulatory comparison may be well-defined even in the absence of shared metric structure. \\
	\textbf{Dependencies:} coherence representation; regulatory ordering; order-preserving map. \\
	\textbf{Assumptions:} ordinal adequacy of representation for regulation; monotone inter-system
	mapping. \\
	\textbf{Main source:} \emph{Coherence Representation in Multi-System Regulation}. \\
	\textbf{Proof status:} P. \\
	\textbf{Note:} This is the main formal proposition grounding multi-system comparability.
	
	\subsection*{R5.06. Multi-system co-regulation operates through aligned regulatory variables rather than through shared global structure}
	\textbf{Type:} Proposition \\
	\textbf{Statement:} When two systems participate in coordinated regulation, the operative basis
	of that coordination need not be a shared global architecture or a shared metric reconstruction
	of coherence. It is sufficient that their internally available regulatory variables admit
	adequate order alignment for coordinated sensitivity to regulatory significance. \\
	\textbf{Dependencies:} inter-system comparability; order-preserving mapping; coherence
	representation. \\
	\textbf{Assumptions:} ordinal adequacy and usable inter-system alignment of regulatory
	variables. \\
	\textbf{Main source:} \emph{Coherence Representation in Multi-System Regulation}. \\
	\textbf{Proof status:} S/C. \\
	\textbf{Note:} This proposition expresses the operational meaning of co-regulation in the
	multi-system branch and is best treated as a canonical normalized formulation.
	
	\subsection*{R5.07. Inter-system conflict is incompatibility of admissible discrepancy domains sufficient for coordinated regulation}
	\textbf{Type:} Proposition/Definitional result \\
	\textbf{Statement:} Inter-system conflict is not defined primarily as disagreement of explicit
	representations, semantic contradiction, or opposition of reported states. It is defined as
	incompatibility between the admissible discrepancy domains of interacting systems such that no
	shared discrepancy structure sufficient for coordinated regulation is available. \\
	\textbf{Dependencies:} admissible discrepancy domain; regulatory accessibility; multi-system
	interaction setting. \\
	\textbf{Assumptions:} conflict evaluated at the level of admissible discrepancy structure rather
	than at the level of explicit report alone. \\
	\textbf{Main source:} \emph{Inter-System Conflict Geometry in Cognitive Systems}. \\
	\textbf{Proof status:} P/S. \\
	\textbf{Note:} This proposition fixes the architectural level of the inter-system conflict
	branch and prevents reduction to representational disagreement.
	
	\subsection*{R5.08. Inter-system conflict may exist without direct representation in one or both systems}
	\textbf{Type:} Proposition \\
	\textbf{Statement:} Because conflict is defined over admissible discrepancy structures rather
	than only over explicit representation, inter-system conflict may exist even when one or both
	systems do not directly represent the conflict as such. In this sense, conflict may be
	structurally operative without being symmetrically represented. \\
	\textbf{Dependencies:} inter-system conflict criterion; restricted accessibility; admissible
	discrepancy domains. \\
	\textbf{Assumptions:} limited accessibility of discrepancy structures within one or more
	participating systems. \\
	\textbf{Main source:} \emph{Inter-System Conflict Geometry in Cognitive Systems}. \\
	\textbf{Proof status:} S/C. \\
	\textbf{Note:} This proposition extends the restricted-accessibility logic into the inter-system
	branch.
	
	\subsection*{R5.09. Alignment and conflict are distinct extension problems}
	\textbf{Type:} Proposition \\
	\textbf{Statement:} Inter-system order alignment and inter-system conflict do not represent two
	descriptions of the same phenomenon. Alignment concerns the existence of a transferable
	regulatory ordering sufficient for coordinated regulation, whereas conflict concerns the absence
	of a shared admissible discrepancy structure sufficient for such coordination. They are
	therefore distinct extension problems of the framework. \\
	\textbf{Dependencies:} inter-system comparability; order alignment; inter-system conflict
	criterion. \\
	\textbf{Assumptions:} none beyond the distinction between regulatory order transfer and
	admissible discrepancy incompatibility. \\
	\textbf{Main source:} \emph{Coherence Representation in Multi-System Regulation} together with
	\emph{Inter-System Conflict Geometry in Cognitive Systems}. \\
	\textbf{Proof status:} C. \\
	\textbf{Note:} This is a canonical collected proposition clarifying the architecture of the
	multi-system branch as a whole.
	
	\section{Collected synthetic framework-level results}
	
	\subsection*{R6.01. The full framework is organized by layered admissibility}
	\textbf{Type:} Proposition / Synthetic result \\
	\textbf{Statement:} The \emph{General Theory of Cognitive Structuring} is not organized by a
	single undifferentiated notion of constraint or updating. Rather, it is organized by layered
	admissibility. At minimum, this layered structure includes:
	\begin{enumerate}[leftmargin=2em]
		\item geometric or structural inconsistency relative to admissible organization,
		\item admissibility of discrepancy for regulation,
		\item representational construction from admitted signals,
		\item admissibility of structural updating.
	\end{enumerate}
	These layers are not interchangeable and should not be collapsed into one operator or one level
	of description. \\
	\textbf{Dependencies:} geometric coherence; pre-symbolic admissibility; admissible discrepancy
	domain; coherence representation; admissibility of structural updating. \\
	\textbf{Assumptions:} cross-series integration of geometry, access, representation, and
	structural updating. \\
	\textbf{Main source:} \emph{General Theory of Cognitive Structuring}. \\
	\textbf{Proof status:} C. \\
	\textbf{Note:} This is the main synthetic proposition of the full framework. It is not a new
	foundational theorem but a normalized convergence claim across earlier branches.
	
	\subsection*{R6.02. Structural inconsistency, discrepancy accessibility, representation, and structural updating are distinct layers}
	\textbf{Type:} Proposition / Synthetic result \\
	\textbf{Statement:} A system may be structurally outside its admissible region without thereby
	having regulatorily accessible discrepancy, may have admitted discrepancy without thereby
	constructing a stable internal representation of that discrepancy, and may possess such a
	representation without structural updating being admissible. Therefore, structural
	inconsistency, discrepancy accessibility, representation, and structural updating are distinct
	layers of the architecture and must not be treated as automatic continuations of one another. \\
	\textbf{Dependencies:} restricted accessibility of coherence; coherence representation;
	pre-symbolic admissibility; admissibility of structural updating. \\
	\textbf{Assumptions:} layered admissibility architecture; non-equivalence of access, representation,
	and updating operators. \\
	\textbf{Main source:} \emph{General Theory of Cognitive Structuring}. \\
	\textbf{Proof status:} C. \\
	\textbf{Note:} This is a synthetic anti-collapse proposition. Its role is to stabilize the
	series against false linearization of the processing chain.
	
	\subsection*{R6.03. Representation is defined over a restricted admissible domain and does not reconstruct full structural deviation}
	\textbf{Type:} Proposition / Synthetic result \\
	\textbf{Statement:} Internal representation of structural pressure is constructed not from the
	full space of structural deviations, but from the restricted domain of discrepancies admitted
	for further processing. Accordingly, representation does not reconstruct the full structure of
	deviation from admissible organization, but only a constrained image of it sufficient for
	regulation under current architectural limits. \\
	\textbf{Dependencies:} admissible discrepancy domain; coherence representation; restricted
	accessibility of coherence. \\
	\textbf{Assumptions:} pre-symbolic admissibility precedes stable representation. \\
	\textbf{Main source:} \emph{General Theory of Cognitive Structuring}. \\
	\textbf{Proof status:} C. \\
	\textbf{Note:} This synthetic proposition unifies the representation and access branches into a
	single canonical constraint statement.
	
	\subsection*{R6.04. Overload introduces historical depth into admissibility across the framework}
	\textbf{Type:} Proposition / Synthetic result \\
	\textbf{Statement:} Because overload memory is retained as a genuine state component,
	admissibility throughout the framework acquires historical depth. The current status of a
	system cannot, in general, be inferred from present structural pressure alone; it depends on
	the retained trace of prior non-compensated load and on the trajectory that produced that
	trace. \\
	\textbf{Dependencies:} overload memory; non-reducibility of overload memory; trajectory-dependent
	regulation; admissibility operator. \\
	\textbf{Assumptions:} overload retention and joint-state admissibility. \\
	\textbf{Main source:} \emph{General Theory of Cognitive Structuring}, grounded in the overload
	and trajectory papers. \\
	\textbf{Proof status:} C. \\
	\textbf{Note:} This is the synthetic version of the historical-state claim running through the
	regulatory branch.
	
	\subsection*{R6.05. Invariant structure is the architectural source of admissible geometry and compression is the mechanism of its non-equivalent transformation}
	\textbf{Type:} Proposition / Synthetic result \\
	\textbf{Statement:} Invariant structure is the architectural source of admissible geometry,
	while structural compression is the principal mechanism through which invariant organization may
	be transformed in a stabilizing way. When such transformation changes admissible geometry
	non-equivalently, the result is level transition rather than merely local structural reduction. \\
	\textbf{Dependencies:} invariant-induced stability geometry; compression operation; compression
	vs simplification; level transition criterion. \\
	\textbf{Assumptions:} invariant-grounded geometry and admissible compression framework. \\
	\textbf{Main source:} \emph{General Theory of Cognitive Structuring}. \\
	\textbf{Proof status:} C. \\
	\textbf{Note:} This proposition synthetically links the structural architecture branch to the
	architectural evolution branch.
	
	\subsection*{R6.06. Identity, representation, and inter-system regulation are downstream consequences of a common layered architecture}
	\textbf{Type:} Proposition / Synthetic result \\
	\textbf{Statement:} Identity dynamics, internal coherence representation, and inter-system
	regulation should not be treated as separate theoretical domains joined only externally. They
	are downstream consequences of a common layered architecture in which admissible geometry,
	overload history, discrepancy accessibility, and constrained representation jointly determine
	the modes through which regulation becomes stable, internally organized, and transferable across
	systems. \\
	\textbf{Dependencies:} layered admissibility architecture; identity as regulatory attractor;
	coherence representation; inter-system order alignment. \\
	\textbf{Assumptions:} synthetic integration across regulatory, representational, and multi-system
	branches. \\
	\textbf{Main source:} \emph{General Theory of Cognitive Structuring}. \\
	\textbf{Proof status:} C. \\
	\textbf{Note:} This is a convergence proposition clarifying that the later branches of the
	series share one architectural backbone rather than merely coexisting as adjacent topics.
	
	\subsection*{R6.07. The General Theory functions as a synthetic convergence framework rather than as the hidden source of earlier branches}
	\textbf{Type:} Meta-structural proposition \\
	\textbf{Statement:} The \emph{General Theory of Cognitive Structuring} does not function as the
	reverse foundational source of the earlier papers in the series. Its role is synthetic: it
	collects and orders earlier results into a layered framework. Accordingly, the dependency
	structure of the series runs from earlier geometric, regulatory, invariant, representational,
	and inter-system branches toward the synthetic convergence paper, not the reverse. \\
	\textbf{Dependencies:} acyclicity criteria; dependency architecture of the series; layered
	integration of prior branches. \\
	\textbf{Assumptions:} none beyond the dependency standards adopted in the verification package. \\
	\textbf{Main source:} companion verification documents together with \emph{General Theory of
		Cognitive Structuring}. \\
	\textbf{Proof status:} C. \\
	\textbf{Note:} This is a meta-structural proposition included to stabilize canonical reading of
	the series order and to prevent false reverse dependency attribution.
	
	\section{Master table of collected propositions and theorems}
	
	\begin{longtable}{>{\RaggedRight\arraybackslash}p{0.10\textwidth}
			>{\RaggedRight\arraybackslash}p{0.31\textwidth}
			>{\RaggedRight\arraybackslash}p{0.13\textwidth}
			>{\RaggedRight\arraybackslash}p{0.08\textwidth}
			>{\RaggedRight\arraybackslash}p{0.17\textwidth}
			>{\RaggedRight\arraybackslash}p{0.14\textwidth}}
		\toprule
		\textbf{Result ID} & \textbf{Title} & \textbf{Type} & \textbf{Status} & \textbf{Main source} & \textbf{Depends on} \\
		\midrule
		\endfirsthead
		\toprule
		\textbf{Result ID} & \textbf{Title} & \textbf{Type} & \textbf{Status} & \textbf{Main source} & \textbf{Depends on} \\
		\midrule
		\endhead
		
		G1.01 & Coherence as distance to the stability region & Definition-formal proposition & P & \emph{Coherence Evaluation in Cognitive Systems} & configuration space; stability region \\
		
		G1.02 & Coherence is distinct from structural tension & Proposition & C & \emph{Coherence Evaluation in Cognitive Systems} & geometric coherence; tension decomposition \\
		
		R2.01 & Structural admissibility depends on the joint regulatory state & Proposition & P & \emph{Structural Admissibility in Cognitive Systems} & tension; overload memory; admissibility operator \\
		
		R2.02 & Overload memory is not reducible to instantaneous tension & Proposition & P & \emph{Overload Formation in Cognitive Processing} & overload excess; overload update rule \\
		
		R2.03 & Regulatory state is trajectory-dependent & Proposition & P & \emph{Trajectory-Dependent Regulation in Cognitive Systems} & overload memory; admissibility operator \\
		
		R2.04 & Overload deforms the admissibility boundary and introduces hysteresis & Proposition & P/S & \emph{Trajectory-Dependent Regulation in Cognitive Systems} & stability region; critical boundary; overload memory \\
		
		R2.05 & The regulatory phase space is minimally two-dimensional & Proposition & C & \emph{Trajectory-Dependent Regulation in Cognitive Systems}, grounded in earlier regulatory papers & tension; overload memory; joint-state admissibility \\
		
		R3.01 & Invariant-induced stability geometry & Proposition /Theorem & P & \emph{Invariants in Cognitive Architectures} & invariants; admissibility profiles; thresholds \\
		
		R3.02 & Global conflict is structured by invariant interaction architecture & Proposition & P & \emph{Invariants in Cognitive Architectures} & invariant structure; interaction graph; incompatibility functions \\
		
		R3.03 & Compression is not equivalent to simplification & Proposition & P & \emph{Structural Compression in Cognitive Architectures} & admissible updating; expected overload criterion \\
		
		R3.04 & Compression generates a new invariant & Proposition & P & \emph{Structural Compression in Cognitive Architectures} & compression operation; invariant definition \\
		
		R3.05 & Compression can preserve or transform admissible geometry & Proposition & P/S & \emph{Structural Compression in Cognitive Architectures} & compression; admissible geometry; architectural level \\
		
		R3.06 & Level transition is a non-equivalent change in admissible geometry & Proposition & P/C & \emph{Structural Admissibility in Cognitive Systems}, strengthened in \emph{Structural Compression in Cognitive Architectures} & architectural level; geometric equivalence \\
		
		R3.07 & Compression improves effective long-run regulatory stability & Proposition & P/S & \emph{Structural Compression in Cognitive Architectures} & compression criterion; expected overload \\
		
		R4.01 & Coherence representation is an internal regulatory variable rather than a reconstruction of geometric coherence & Proposition & P & \emph{Emergence of Coherence Representation} & geometric coherence; observable signals; representation map \\
		
		R4.02 & Coherence representation preserves regulatory order without requiring metric preservation & Proposition & P/C & \emph{Emergence of Coherence Representation}, strengthened in \emph{Coherence Representation in Multi-System Regulation} & coherence representation; regulatory ordering \\
		
		R4.03 & Pre-symbolic admissibility is a distinct prior filtering layer & Proposition-like definitional result & D/C & \emph{Pre-Symbolic Admissibility in Cognitive Systems} & signal domain; regulatory layering \\
		
		R4.04 & Pre-symbolic admissibility determines the admissible discrepancy domain & Proposition & P/S & \emph{Pre-Symbolic Admissibility in Cognitive Systems} & pre-symbolic signal domain; pre-symbolic operator \\
		
		R4.05 & Restricted accessibility of coherence & Proposition & P & \emph{Restricted Accessibility of Coherence in Cognitive Systems} & geometric coherence; admissible discrepancy domain; regulatory accessibility \\
		
		R4.06 & Non-injectivity of regulatory accessibility & Proposition & P/S & \emph{Restricted Accessibility of Coherence in Cognitive Systems} & restricted accessibility; admissible discrepancy structure \\
		
		R4.07 & Restricted accessibility is structural, not merely representationally approximate & Proposition & C & \emph{Restricted Accessibility of Coherence in Cognitive Systems} & restricted accessibility; non-injectivity; pre-symbolic admissibility \\
		
		R4.08 & Repeated admissible propagation induces propagation bias & Proposition & S & \emph{Pre-Symbolic Admissibility in Cognitive Systems} & pre-symbolic admissibility; repeated propagation \\
		
		R4.09 & Repeated admissible propagation may stabilize proto-structural channels & Proposition & S/O & \emph{Pre-Symbolic Admissibility in Cognitive Systems} & propagation bias; repeated admissible selection \\
		
		R5.01 & Identity is derived as a regulatory attractor rather than introduced as a primitive self-representational object & Proposition & S/C & \emph{Identity as a Regulatory Attractor} & overload; trajectory dynamics; bounded regulation \\
		
		R5.02 & Low-overload concentration regions may function as identity-cores & Proposition /Theorem & S/P & \emph{Identity as a Regulatory Attractor} & overload dynamics; admissibility structure; trajectory dependence \\
		
		R5.03 & Identity persistence is regulatory rather than merely representational & Proposition & C & \emph{Identity as a Regulatory Attractor} & identity as attractor; low-overload concentration regions \\
		
		R5.04 & Inter-system comparability does not require shared metric coherence or shared configuration space & Proposition & P/C & \emph{Coherence Representation in Multi-System Regulation} & coherence representation; non-metric ordering \\
		
		R5.05 & Minimal transferable structure between systems is ordinal rather than metric & Proposition & P & \emph{Coherence Representation in Multi-System Regulation} & coherence representation; order-preserving map \\
		
		R5.06 & Multi-system co-regulation operates through aligned regulatory variables rather than through shared global structure & Proposition & S/C & \emph{Coherence Representation in Multi-System Regulation} & inter-system comparability; order alignment \\
		
		R5.07 & Inter-system conflict is incompatibility of admissible discrepancy domains sufficient for coordinated regulation & Proposition /Definitional result & P/S & \emph{Inter-System Conflict Geometry in Cognitive Systems} & admissible discrepancy domain; multi-system interaction \\
		
		R5.08 & Inter-system conflict may exist without direct representation in one or both systems & Proposition & S/C & \emph{Inter-System Conflict Geometry in Cognitive Systems} & inter-system conflict criterion; restricted accessibility \\
		
		R5.09 & Alignment and conflict are distinct extension problems & Proposition & C & \emph{Coherence Representation in Multi-System Regulation} together with \emph{Inter-System Conflict Geometry in Cognitive Systems} & order alignment; inter-system conflict criterion \\
		
		R6.01 & The full framework is organized by layered admissibility & Proposition / Synthetic result & C & \emph{General Theory of Cognitive Structuring} & geometry; access; representation; structural updating \\
		
		R6.02 & Structural inconsistency, discrepancy accessibility, representation, and structural updating are distinct layers & Proposition / Synthetic result & C & \emph{General Theory of Cognitive Structuring} & restricted accessibility; representation; updating admissibility \\
		
		R6.03 & Representation is defined over a restricted admissible domain and does not reconstruct full structural deviation & Proposition / Synthetic result & C & \emph{General Theory of Cognitive Structuring} & admissible discrepancy domain; representation; restricted accessibility \\
		
		R6.04 & Overload introduces historical depth into admissibility across the framework & Proposition / Synthetic result & C & \emph{General Theory of Cognitive Structuring}, grounded in overload and trajectory papers & overload memory; trajectory dependence; admissibility operator \\
		
		R6.05 & Invariant structure is the architectural source of admissible geometry and compression is the mechanism of its non-equivalent transformation & Proposition / Synthetic result & C & \emph{General Theory of Cognitive Structuring} & invariant-induced geometry; compression; level transition \\
		
		R6.06 & Identity, representation, and inter-system regulation are downstream consequences of a common layered architecture & Proposition / Synthetic result & C & \emph{General Theory of Cognitive Structuring} & layered admissibility; identity; representation; inter-system alignment \\
		
		R6.07 & The General Theory functions as a synthetic convergence framework rather than as the hidden source of earlier branches & Meta-structural proposition & C & companion verification documents together with \emph{General Theory of Cognitive Structuring} & dependency architecture; synthetic integration \\
		
		\bottomrule
	\end{longtable}
	
	\section{Proof status note by family}
	
	\subsection{Results with full proof in source}
	
	The following results are currently treated as proof-complete in their source papers and are
	therefore suitable for direct canonical reference, even where wording may later be normalized
	further:
	
	\begin{itemize}[leftmargin=2em]
		\item G1.01 --- Coherence as distance to the stability region
		\item R2.01 --- Structural admissibility depends on the joint regulatory state
		\item R2.02 --- Overload memory is not reducible to instantaneous tension
		\item R2.03 --- Regulatory state is trajectory-dependent
		\item R3.01 --- Invariant-induced stability geometry
		\item R3.02 --- Global conflict is structured by invariant interaction architecture
		\item R3.03 --- Compression is not equivalent to simplification
		\item R3.04 --- Compression generates a new invariant
		\item R4.01 --- Coherence representation is an internal regulatory variable rather than a reconstruction of geometric coherence
		\item R4.05 --- Restricted accessibility of coherence
		\item R5.05 --- Minimal transferable structure between systems is ordinal rather than metric
	\end{itemize}
	
	\subsection{Results stated with proof sketch or partially compressed support}
	
	The following results are supported in their source papers but remain, in their current collected
	form, sketch-level, assumption-local, or in need of cleaner canonical normalization:
	
	\begin{itemize}[leftmargin=2em]
		\item R2.04 --- Overload deforms the admissibility boundary and introduces hysteresis
		\item R3.05 --- Compression can preserve or transform admissible geometry
		\item R3.07 --- Compression improves effective long-run regulatory stability
		\item R4.04 --- Pre-symbolic admissibility determines the admissible discrepancy domain
		\item R4.06 --- Non-injectivity of regulatory accessibility
		\item R4.08 --- Repeated admissible propagation induces propagation bias
		\item R4.09 --- Repeated admissible propagation may stabilize proto-structural channels
		\item R5.01 --- Identity is derived as a regulatory attractor rather than introduced as a primitive self-representational object
		\item R5.02 --- Low-overload concentration regions may function as identity-cores
		\item R5.06 --- Multi-system co-regulation operates through aligned regulatory variables rather than through shared global structure
		\item R5.07 --- Inter-system conflict is incompatibility of admissible discrepancy domains sufficient for coordinated regulation
		\item R5.08 --- Inter-system conflict may exist without direct representation in one or both systems
	\end{itemize}
	
	\subsection{Results canonically normalized here from distributed support}
	
	The following results are collected in canonical form from one or more papers. Their role in the
	framework is stable, but the present note should not be read as upgrading them beyond the status
	supported by their sources:
	
	\begin{itemize}[leftmargin=2em]
		\item G1.02 --- Coherence is distinct from structural tension
		\item R2.05 --- The regulatory phase space is minimally two-dimensional
		\item R3.06 --- Level transition is a non-equivalent change in admissible geometry
		\item R4.02 --- Coherence representation preserves regulatory order without requiring metric preservation
		\item R4.07 --- Restricted accessibility is structural, not merely representationally approximate
		\item R5.03 --- Identity persistence is regulatory rather than merely representational
		\item R5.04 --- Inter-system comparability does not require shared metric coherence or shared configuration space
		\item R5.09 --- Alignment and conflict are distinct extension problems
		\item R6.01 --- The full framework is organized by layered admissibility
		\item R6.02 --- Structural inconsistency, discrepancy accessibility, representation, and structural updating are distinct layers
		\item R6.03 --- Representation is defined over a restricted admissible domain and does not reconstruct full structural deviation
		\item R6.04 --- Overload introduces historical depth into admissibility across the framework
		\item R6.05 --- Invariant structure is the architectural source of admissible geometry and compression is the mechanism of its non-equivalent transformation
		\item R6.06 --- Identity, representation, and inter-system regulation are downstream consequences of a common layered architecture
		\item R6.07 --- The General Theory functions as a synthetic convergence framework rather than as the hidden source of earlier branches
	\end{itemize}
	
	\subsection{Definition-like or mixed-status inclusions kept for structural completeness}
	
	The following items are included in the collected propositions and theorems layer because they
	function as structurally central separating results, even though they retain a definitional or
	mixed status:
	
	\begin{itemize}[leftmargin=2em]
		\item R4.03 --- Pre-symbolic admissibility is a distinct prior filtering layer
	\end{itemize}
	
	\subsection{Results most in need of later proof-compendium consolidation}
	
	The following results are likely to benefit most from later extraction into a full normalized
	proof layer:
	
	\begin{itemize}[leftmargin=2em]
		\item R2.04 --- Overload deforms the admissibility boundary and introduces hysteresis
		\item R3.05 --- Compression can preserve or transform admissible geometry
		\item R3.07 --- Compression improves effective long-run regulatory stability
		\item R4.06 --- Non-injectivity of regulatory accessibility
		\item R4.08 --- Repeated admissible propagation induces propagation bias
		\item R4.09 --- Repeated admissible propagation may stabilize proto-structural channels
		\item R5.01 --- Identity is derived as a regulatory attractor rather than introduced as a primitive self-representational object
		\item R5.02 --- Low-overload concentration regions may function as identity-cores
		\item R5.06 --- Multi-system co-regulation operates through aligned regulatory variables rather than through shared global structure
		\item R5.07 --- Inter-system conflict is incompatibility of admissible discrepancy domains sufficient for coordinated regulation
		\item R5.08 --- Inter-system conflict may exist without direct representation in one or both systems
	\end{itemize}
	
	\section{Results reserved for future proof compendium}
	
	The following result families are the most natural candidates for later extraction into a full
	proof compendium.
	
	\subsection*{1. Core geometric and regulatory results}
	
	These results form the formal backbone of the framework and should be among the first to receive
	full normalized proof treatment:
	
	\begin{itemize}[leftmargin=2em]
		\item G1.01 --- Coherence as distance to the stability region
		\item G1.02 --- Coherence is distinct from structural tension
		\item R2.01 --- Structural admissibility depends on the joint regulatory state
		\item R2.02 --- Overload memory is not reducible to instantaneous tension
		\item R2.03 --- Regulatory state is trajectory-dependent
		\item R2.04 --- Overload deforms the admissibility boundary and introduces hysteresis
		\item R2.05 --- The regulatory phase space is minimally two-dimensional
	\end{itemize}
	
	\subsection*{2. Invariant and compression results}
	
	These results form the architectural transformation branch and are likely to require a compact
	lemma structure in a later proof layer:
	
	\begin{itemize}[leftmargin=2em]
		\item R3.01 --- Invariant-induced stability geometry
		\item R3.02 --- Global conflict is structured by invariant interaction architecture
		\item R3.03 --- Compression is not equivalent to simplification
		\item R3.04 --- Compression generates a new invariant
		\item R3.05 --- Compression can preserve or transform admissible geometry
		\item R3.06 --- Level transition is a non-equivalent change in admissible geometry
		\item R3.07 --- Compression improves effective long-run regulatory stability
	\end{itemize}
	
	\subsection*{3. Representation and accessibility results}
	
	These results are especially important because they fix the relation between geometry, access,
	and internal regulatory approximation:
	
	\begin{itemize}[leftmargin=2em]
		\item R4.01 --- Coherence representation is an internal regulatory variable rather than a reconstruction of geometric coherence
		\item R4.02 --- Coherence representation preserves regulatory order without requiring metric preservation
		\item R4.04 --- Pre-symbolic admissibility determines the admissible discrepancy domain
		\item R4.05 --- Restricted accessibility of coherence
		\item R4.06 --- Non-injectivity of regulatory accessibility
		\item R4.07 --- Restricted accessibility is structural, not merely representationally approximate
		\item R4.08 --- Repeated admissible propagation induces propagation bias
		\item R4.09 --- Repeated admissible propagation may stabilize proto-structural channels
	\end{itemize}
	
	\subsection*{4. Identity and multi-system results}
	
	These results are likely to require the most careful later normalization because they involve
	either long-run assumptions or cross-system transfer conditions:
	
	\begin{itemize}[leftmargin=2em]
		\item R5.01 --- Identity is derived as a regulatory attractor rather than introduced as a primitive self-representational object
		\item R5.02 --- Low-overload concentration regions may function as identity-cores
		\item R5.03 --- Identity persistence is regulatory rather than merely representational
		\item R5.04 --- Inter-system comparability does not require shared metric coherence or shared configuration space
		\item R5.05 --- Minimal transferable structure between systems is ordinal rather than metric
		\item R5.06 --- Multi-system co-regulation operates through aligned regulatory variables rather than through shared global structure
		\item R5.07 --- Inter-system conflict is incompatibility of admissible discrepancy domains sufficient for coordinated regulation
		\item R5.08 --- Inter-system conflict may exist without direct representation in one or both systems
		\item R5.09 --- Alignment and conflict are distinct extension problems
	\end{itemize}
	
	\subsection*{5. Synthetic framework-level results}
	
	These results are best treated as canonical convergence statements and may or may not belong to a
	future proof compendium in theorem-grade form. Their inclusion should remain explicitly
	non-foundational:
	
	\begin{itemize}[leftmargin=2em]
		\item R6.01 --- The full framework is organized by layered admissibility
		\item R6.02 --- Structural inconsistency, discrepancy accessibility, representation, and structural updating are distinct layers
		\item R6.03 --- Representation is defined over a restricted admissible domain and does not reconstruct full structural deviation
		\item R6.04 --- Overload introduces historical depth into admissibility across the framework
		\item R6.05 --- Invariant structure is the architectural source of admissible geometry and compression is the mechanism of its non-equivalent transformation
		\item R6.06 --- Identity, representation, and inter-system regulation are downstream consequences of a common layered architecture
		\item R6.07 --- The General Theory functions as a synthetic convergence framework rather than as the hidden source of earlier branches
	\end{itemize}
	
	\section{Compact reference table}
	
	\begin{longtable}{>{\RaggedRight\arraybackslash}p{0.11\textwidth}
			>{\RaggedRight\arraybackslash}p{0.47\textwidth}
			>{\RaggedRight\arraybackslash}p{0.10\textwidth}
			>{\RaggedRight\arraybackslash}p{0.18\textwidth}}
		\toprule
		\textbf{ID} & \textbf{Title} & \textbf{Status} & \textbf{Source} \\
		\midrule
		\endfirsthead
		\toprule
		\textbf{ID} & \textbf{Title} & \textbf{Status} & \textbf{Source} \\
		\midrule
		\endhead
		
		G1.01 & Coherence as distance to the stability region & P & coherence paper \\
		
		G1.02 & Coherence is distinct from structural tension & C & coherence paper \\
		
		R2.01 & Structural admissibility depends on the joint regulatory state & P & structural admissibility \\
		
		R2.02 & Overload memory is not reducible to instantaneous tension & P & overload formation \\
		
		R2.03 & Regulatory state is trajectory-dependent & P & trajectory regulation \\
		
		R2.04 & Overload deforms the admissibility boundary and introduces hysteresis & P/S & trajectory regulation \\
		
		R2.05 & The regulatory phase space is minimally two-dimensional & C & trajectory regulation + earlier regulatory papers \\
		
		R3.01 & Invariant-induced stability geometry & P & invariants paper \\
		
		R3.02 & Global conflict is structured by invariant interaction architecture & P & invariants paper \\
		
		R3.03 & Compression is not equivalent to simplification & P & compression paper \\
		
		R3.04 & Compression generates a new invariant & P & compression paper \\
		
		R3.05 & Compression can preserve or transform admissible geometry & P/S & compression paper \\
		
		R3.06 & Level transition is a non-equivalent change in admissible geometry & P/C & structural admissibility + compression \\
		
		R3.07 & Compression improves effective long-run regulatory stability & P/S & compression paper \\
		
		R4.01 & Coherence representation is an internal regulatory variable rather than a reconstruction of geometric coherence & P & emergence paper \\
		
		R4.02 & Coherence representation preserves regulatory order without requiring metric preservation & P/C & emergence paper + multi-system regulation \\
		
		R4.03 & Pre-symbolic admissibility is a distinct prior filtering layer & D/C & pre-symbolic paper \\
		
		R4.04 & Pre-symbolic admissibility determines the admissible discrepancy domain & P/S & pre-symbolic paper \\
		
		R4.05 & Restricted accessibility of coherence & P & restricted accessibility paper \\
		
		R4.06 & Non-injectivity of regulatory accessibility & P/S & restricted accessibility paper \\
		
		R4.07 & Restricted accessibility is structural, not merely representationally approximate & C & restricted accessibility paper \\
		
		R4.08 & Repeated admissible propagation induces propagation bias & S & pre-symbolic paper \\
		
		R4.09 & Repeated admissible propagation may stabilize proto-structural channels & S/O & pre-symbolic paper \\
		
		R5.01 & Identity is derived as a regulatory attractor rather than introduced as a primitive self-representational object & S/C & identity paper \\
		
		R5.02 & Low-overload concentration regions may function as identity-cores & S/P & identity paper \\
		
		R5.03 & Identity persistence is regulatory rather than merely representational & C & identity paper \\
		
		R5.04 & Inter-system comparability does not require shared metric coherence or shared configuration space & P/C & multi-system regulation \\
		
		R5.05 & Minimal transferable structure between systems is ordinal rather than metric & P & multi-system regulation \\
		
		R5.06 & Multi-system co-regulation operates through aligned regulatory variables rather than through shared global structure & S/C & multi-system regulation \\
		
		R5.07 & Inter-system conflict is incompatibility of admissible discrepancy domains sufficient for coordinated regulation & P/S & inter-system conflict \\
		
		R5.08 & Inter-system conflict may exist without direct representation in one or both systems & S/C & inter-system conflict \\
		
		R5.09 & Alignment and conflict are distinct extension problems & C & multi-system regulation + inter-system conflict \\
		
		R6.01 & The full framework is organized by layered admissibility & C & general theory \\
		
		R6.02 & Structural inconsistency, discrepancy accessibility, representation, and structural updating are distinct layers & C & general theory \\
		
		R6.03 & Representation is defined over a restricted admissible domain and does not reconstruct full structural deviation & C & general theory \\
		
		R6.04 & Overload introduces historical depth into admissibility across the framework & C & general theory + overload / trajectory papers \\
		
		R6.05 & Invariant structure is the architectural source of admissible geometry and compression is the mechanism of its non-equivalent transformation & C & general theory \\
		
		R6.06 & Identity, representation, and inter-system regulation are downstream consequences of a common layered architecture & C & general theory \\
		
		R6.07 & The General Theory functions as a synthetic convergence framework rather than as the hidden source of earlier branches & C & verification package + general theory \\
		
		\bottomrule
	\end{longtable}
	
	
	
\end{document}